\chapter[Referencial Teórico]{Referencial Teórico}

Este capítulo apresenta a metodologia utilizada para conduzir o desenvolvimento deste trabalho de conclusão de curso. Inicialmente, descreve-se a classificação da pesquisa e as etapas que compõem o planejamento e execução do projeto. Em seguida, detalham-se os procedimentos adotados para a construção do referencial teórico, bem como a metodologia aplicada no desenvolvimento do sistema proposto, que segue uma abordagem ágil visando maior flexibilidade e eficiência. Por fim, aborda-se a estratégia de gerência de configuração e evolução de software, essencial para assegurar a qualidade, rastreabilidade e facilidade de manutenção do sistema ao longo do tempo.

\section{Login Único}

O Login Único, iniciativa desenvolvida pelo Governo Federal, apresenta-se como uma solução estratégica de autenticação digital voltada à centralização do acesso do cidadão aos serviços públicos eletrônicos. Atuando como uma credencial universal, essa plataforma visa simplificar a experiência do usuário, promovendo a integração entre os diversos sistemas governamentais. 

De acordo com o Roteiro Técnico de Solicitação de Credencial de Processo, a integração de sistemas ao Login Único segue um fluxo estruturado que inclui solicitação formal, validação em ambiente de homologação, envio de evidências técnicas e aprovação para uso em produção. Essa estrutura não apenas garante a conformidade técnica dos serviços integrados, como também assegura a proteção dos dados pessoais, em consonância com os princípios da segurança da informação. A arquitetura do Login Único baseia-se no protocolo OAuth 2.0, amplamente adotado no contexto da autenticação digital moderna, permitindo a implementação de fluxos seguros de autorização e gerenciamento de sessões. A plataforma ainda prevê níveis diferenciados de identidade (bronze, prata e ouro), os quais definem a profundidade do acesso conforme o grau de confiabilidade atribuído ao usuário.

Nesse sentido, o Login Único não apenas atua como facilitador do acesso, mas também como um componente fundamental da infraestrutura de cidadania digital no Brasil, promovendo eficiência, padronização e confiança no relacionamento entre Estado e sociedade.~\cite{govbrlogin}


\section{Assinatura Digital}

A assinatura digital é uma tecnologia utilizada para garantir a autenticidade, integridade e não-repúdio de documentos e transações realizadas no meio eletrônico. Ela utiliza a criptografia assimétrica, onde um par de chaves é utilizado: uma chave privada, que é mantida em sigilo pelo proprietário, e uma chave pública, que é disponibilizada para verificação. A assinatura digital é gerada a partir da chave privada do assinante e serve para garantir que o documento ou transação não foi alterado desde que foi assinado, além de confirmar a identidade do signatário.

Essa tecnologia é amplamente aplicada em diversos setores, como jurídico, financeiro e governamental, pois oferece uma maneira segura e confiável de realizar transações sem a necessidade de documentos físicos. No Brasil, a \textit{Infraestrutura de Chaves Públicas Brasileira (ICP-Brasil)}, estabelecida pela Medida Provisória 2.200-2, é a base para a utilização de certificados digitais, que permitem a assinatura de documentos eletrônicos de forma válida e legal, equivalente à assinatura manuscrita \cite{souza2020}.

A assinatura digital é fundamental para a digitalização de processos e para garantir a segurança e a integridade das informações trocadas. Ela também permite que o usuário tenha o controle sobre suas ações e documentos, assegurando que apenas ele possa assinar ou autorizar determinadas transações. Em sistemas jurídicos, a assinatura digital tem validade legal, de acordo com a legislação brasileira, e é frequentemente usada para autenticar contratos, declarações fiscais, e outros documentos importantes.

\subsection{Assinatura Digital no Governo Federal}

O Governo Federal do Brasil adotou a \textit{Assinatura Digital} como uma das principais ferramentas para garantir a autenticidade de documentos e a segurança de transações no âmbito público. Uma das principais plataformas para a implementação da assinatura digital nos serviços governamentais é a \textit{API de Assinatura Digital do GOV.BR}, que integra o \textit{Login Único} e permite a assinatura de documentos eletrônicos por cidadãos e servidores públicos.

Essa plataforma utiliza os padrões de segurança estabelecidos pela ICP-Brasil, permitindo que documentos sejam assinados digitalmente por meio de certificados emitidos por Autoridades Certificadoras credenciadas. O uso dessa tecnologia visa facilitar o processo de digitalização de serviços públicos e promover a eficiência na prestação de serviços ao cidadão. A assinatura digital proporcionada pelo GOV.BR garante não apenas a autenticidade dos documentos, mas também assegura que os dados sensíveis, como informações pessoais, estejam protegidos contra fraudes e acessos não autorizados.

De acordo com o ~\cite{govbrassinaeletronica}, "a plataforma de Assinatura Eletrônica para órgãos oferece uma solução eficaz e segura para a autenticação de documentos, permitindo que os órgãos públicos assinem digitalmente documentos com validade jurídica, conforme os requisitos da ICP-Brasil".